\SourceLeaf{33} from Mr \textsc{Mehmke} informed me that \textsc{Gauss}, in the
second table (apart from the three cases mentioned), also gave for the system
of units representable with integer $a$, $b$, $c$
in the form $a+b\nu+c\nu\nu$ the unit
reciprocal to the fundamental unit in each case. In most cases the unit in question is also the
fundamental unit of the field.
Exceptions to this rule occur at $n=16$, where the eighth power of the
fundamental unit is given, and at $n=10$ and $n=19$, for which the
reciprocal fundamental units are as follows.

\[
\frac{7+\sqrt[3]{10}-2\sqrt[3]{100}}{3}, \qquad
\frac{2+2\sqrt[3]{19}-\sqrt[3]{19^2}}{3}.
\]

\par\noindent\hfill\textsc{Fricke.}\par\endgroup
